\chapter[Band 10: Natürliche Zahlen]{Band 10 auf einen Blick}
\noindent\textbf{Natürliche Zahlen}\par\medskip
\noindent\emph{Lesart.} Der Band unterscheidet die mengentheoretische
Konstruktion und den anschließenden Aufbau mit der Peano-Schnittstelle.
In den arithmetischen Zeilen sind \(a,b,c,m,n\in\mathbb N\);
\(u\colon\mathbb N\to\mathbb N\) ist eine Folge natürlicher Zahlen.
\begin{BandUebersicht}
\textbf{Induktive Menge und von-Neumann-Konstruktion}
\UebFormel{\begin{aligned}
&\Induktiv(A):\quad\varnothing\in A,\\
&\forall x\in A\,(x\cup\{x\}\in A),\\
&\mathbb N\coloneq\bigcap\{A\mid\Induktiv(A)\},\\
&0\coloneqq\varnothing.
\end{aligned}}
\UebQuellen{\FormulaRefAuto{InductiveSetDef}[def];
\FormulaRefAuto{\mathbb{N} \coloneq \bigcap \{ A \mid \Induktiv(A) \}}[def];
\FormulaRefAuto{0\coloneqq\varnothing}[def]}
&
\textit{Kleinste induktive Menge}
\UebFormel{\begin{aligned}
&\Induktiv(\mathbb N),\\
&\Induktiv(A)\vdash\mathbb N\subseteq A.
\end{aligned}}
Das Unendlichkeitsaxiom sichert, dass induktive Mengen existieren;
die Schnittkonstruktion liefert ihre gemeinsame kleinste Struktur.
\UebQuellen{\FormulaRefAuto{\Induktiv(\mathbb{N})}[thm];
\FormulaRefAuto{\Induktiv(A) \vdash \mathbb{N} \subseteq A}[thm]}\\
\addlinespace[.8ex]
\textbf{Peano-Schnittstelle}
\(0\in\mathbb N\), \(\Succ\colon\mathbb N\to\mathbb N\),
kein Nachfolger ist null, und \(\Succ\) ist injektiv. Dazu tritt Induktion:
\UebFormel{\begin{aligned}
&P(0)\dsep\forall n\in\mathbb N\,(P(n)\rightarrow P(\Succ(n)))\\
&\hspace{10mm}\vdash\forall n\in\mathbb N\,P(n).
\end{aligned}}
\UebQuellen{\FormulaRefAuto{Peano-Axiome}[def];
\FormulaRefAuto{PeanoInductionAxiom}[ax]}
&
\textit{Eindeutiger Vorgänger jeder Nichtnullzahl}
\UebFormel{\begin{aligned}
&n\neq0\vdash\exists!k\in\mathbb N\,(n=\Succ(k)),\\
&\Succ\colon\mathbb N\bij\mathbb N\setminus\{0\}.
\end{aligned}}
\UebQuellen{\FormulaRefAuto{PeanoNonzeroUniquePredecessor}[thm];
\FormulaRefAuto{PeanoSuccNonzeroBijection}[thm]}\\
\addlinespace[.8ex]
\textbf{Rekursionskern und Rekursionsabbildung}
Für \(m_0\in M\), \(f\colon M\to M\) wird der kleinste
Graphkandidat mit Startpaar \((0,m_0)\) und Abschluss unter
\((n,a)\mapsto(\Succ(n),f(a))\) konstruiert.
Die daraus gewonnene Funktion heißt \(\RecFun{m_0}{f}\).
\UebQuellen{\FormulaRefAuto{RecFunDef}[def]}
&
\textit{Dedekindscher Rekursionssatz}
\UebFormel{\begin{aligned}
&m_0\in M\dsep f\colon M\to M\vdash\\
&\exists!G\colon\mathbb N\to M\,\bigl(G(0)=m_0\\
&\quad\land\forall n\in\mathbb N\,
G(\Succ(n))=f(G(n))\bigr).
\end{aligned}}
\UebQuellen{\FormulaRefAuto{DedekindRecursionTheorem}[thm]}\\
\addlinespace[.8ex]
\textbf{Addition durch Rekursion}
\UebFormel{a+b\coloneqq\RecFun{a}{\Succ}(b).}
\UebQuellen{\FormulaRefAuto{Addition auf \(\mathbb{N}\)}[def]}
&
\textit{Assoziativität und Kommutativität}
\UebFormel{\begin{aligned}
&(a+b)+c=a+(b+c),\\
&a+b=b+a.
\end{aligned}}
Die rekursive Definition führt damit zu den üblichen Rechengesetzen.
\UebQuellen{\FormulaRefAuto{PeanoAddAssociative}[thm];
\FormulaRefAuto{PeanoAddCommutative}[thm]}\\
\addlinespace[.8ex]
\textbf{Nachfolger- und Vorgängerschreibweise}
\(1\) bezeichnet \(\Succ(0)\); \(\Pred=\Succ^{-1}\) wird auf
\(\mathbb N\setminus\{0\}\) definiert. Für \(n\neq0\) ist
\(n-1\) die entsprechende Vorgängerschreibweise.
\UebQuellen{\FormulaRefAuto{PeanoPred}[def];
\FormulaRefAuto{PeanoMinusOneNotation}[def]}
&
\textit{Induktion mit dem Schritt \(n\mapsto n+1\)}
\UebFormel{\begin{aligned}
&P(0)\dsep\forall n\in\mathbb N\,(P(n)\rightarrow P(n+1))\\
&\hspace{12mm}\vdash\forall n\in\mathbb N\,P(n),\\
&n\neq0\vdash(n-1)+1=n.
\end{aligned}}
\UebQuellen{\FormulaRefAuto{PeanoInductionAddOne}[thm];
\FormulaRefAuto{PeanoMinusOneAddOne}[thm]}\\
\addlinespace[.8ex]
\textbf{Natürliche Ordnung}
Der Band führt für die von-Neumann-Zahlen ein:
\UebFormel{\begin{aligned}
m<_{\mathbb N}n&\coloneqq m\in n,\\
m\leq_{\mathbb N}n&\coloneqq m\subseteq n.
\end{aligned}}
Im arithmetischen Aufbau werden die additiven Charakterisierungen
und konkreten Ordnungsgesetze ausgearbeitet.
\UebQuellen{\FormulaRefAuto{NaturalStrictOrderOperator}[def];
\FormulaRefAuto{NaturalOrderOperator}[def]}
&
\textit{Trichotomie und Wohlordnungsprinzip}
\UebFormel{\begin{aligned}
&m<n\lor m=n\lor n<m,\\
&B\subseteq\mathbb N\dsep B\neq\varnothing\\
&\hspace{8mm}\vdash\exists a\in B\,\forall b\in B\,(a\leq b).
\end{aligned}}
\UebQuellen{\FormulaRefAuto{PeanoOrderTrichotomy}[thm];
\FormulaRefAuto{NaturalWellOrderingPrinciple}[thm]}\\
\addlinespace[.8ex]
\textbf{Anfangsabschnitte}
\UebFormel{\NatSeg n\coloneqq\{i\in\mathbb N\mid i<n\}.}
\UebQuellen{\FormulaRefAuto{PeanoNatSeg}[def]}
&
\textit{Endliche Abschnitte fassen nicht ganz \(\mathbb N\)}
\UebFormel{\begin{aligned}
&i\in\mathbb N\vdash
i\in\NatSeg n\leftrightarrow i<n,\\
&\neg\exists F\colon\mathbb N\inj\NatSeg n.
\end{aligned}}
\UebQuellen{\FormulaRefAuto{PeanoNatSegMembership}[thm];
\FormulaRefAuto{PeanoNatSegNoNatInjection}[thm]}\\
\addlinespace[.8ex]
\textbf{Multiplikation durch wiederholte Addition}
\UebFormel{a\cdot b\coloneqq\RecFun{0}{\AddStep{a}}(b).}
Dabei ist \(\AddStep{a}\) der Additionsschritt mit festem Summanden \(a\).
\UebQuellen{\FormulaRefAuto{Multiplikation auf \(\mathbb{N}\)}[def]}
&
\textit{Multiplikative und distributive Gesetze}
\UebFormel{\begin{aligned}
&a\cdot b=b\cdot a,\\
&(a\cdot b)\cdot c=a\cdot(b\cdot c),\\
&a\cdot(b+c)=a\cdot b+a\cdot c.
\end{aligned}}
\UebQuellen{\FormulaRefAuto{PeanoMulCommutative}[thm];
\FormulaRefAuto{PeanoMulAssociative}[thm];
\FormulaRefAuto{PeanoMulLeftDistributive}[thm]}\\
\addlinespace[.8ex]
\textbf{Endliche Summen und Produkte}
\(R_u\) und \(P_u\) akkumulieren die gelesenen Werte durch
Addition beziehungsweise Multiplikation. Definiert wird:
\UebFormel{\begin{aligned}
\sum_{i<n}u(i)&\coloneqq\pi_2(R_u(n)),\\
\prod_{i<n}u(i)&\coloneqq\pi_2(P_u(n)).
\end{aligned}}
\UebQuellen{\FormulaRefAuto{Summenzeichen}[def];
\FormulaRefAuto{Produktzeichen}[def]}
&
\textit{Rekursive Auswertung der Summe}
\UebFormel{\begin{aligned}
&\sum_{i<0}u(i)=0,\\
&\sum_{i<n+1}u(i)=\left(\sum_{i<n}u(i)\right)+u(n).
\end{aligned}}
\UebQuellen{\FormulaRefAuto{u\colon\mathbb{N}\to\mathbb{N}\vdash \sum_{i<0}u(i)=0}[thm];
\FormulaRefAuto{u\colon\mathbb{N}\to\mathbb{N}\dsep n\in\mathbb{N}\vdash \sum_{i<n+1}u(i)=\left(\sum_{i<n}u(i)\right) + u(n)}[thm]}\\
\addlinespace[.8ex]
\textbf{Potenzen und Fakultäten}
\UebFormel{\begin{aligned}
b^n&\coloneqq\prod_{i<n}\ConstFun{\mathbb N}{\mathbb N}{b}(i),\\
n!&\coloneqq\prod_{i<n}(i+1).
\end{aligned}}
\UebQuellen{\FormulaRefAuto{Potenz natürlicher Zahlen}[def];
\FormulaRefAuto{FactorialFunction}[def]}
&
\textit{Null- und Schrittregeln}
\UebFormel{\begin{aligned}
&b^0=1,\qquad b^{n+1}=b^n\cdot b,\\
&0!=1,\qquad(n+1)!=n!\cdot(n+1).
\end{aligned}}
\UebQuellen{\FormulaRefAuto{PeanoPowerZero}[thm];
\FormulaRefAuto{PeanoPowerAddOne}[thm];
\FormulaRefAuto{FactorialRecursion}[thm]}\\
\addlinespace[.8ex]
\textbf{Stellenwertsysteme}
\UebFormel{\begin{aligned}
\BaseOK{b}&\coloneqq b\in\mathbb N\land2\leq b,\\
\DigitSet{b}&\coloneqq\{d\in\mathbb N\mid d<b\},\\
\PlaceTermSeq{b}{a}(i)&\coloneqq a(i)\cdot b^i.
\end{aligned}}
Zahlzeichen werden durch endliche Summen dieser Stellenwerte
ausgewertet; Dual- und Dezimalsystem sind die Fälle \(b=2,10\).
\UebQuellen{\FormulaRefAuto{Basis eines Stellenwertsystems}[def];
\FormulaRefAuto{Ziffernmenge zur Basis \(b\)}[def];
\FormulaRefAuto{Funktionsvorschrift der Stellenwertfolge}[def]}
&
\textit{Typisierung der Stellenwertfolge}
\UebFormel{\begin{aligned}
&\BaseOK{b}\dsep a\colon\mathbb N\to\mathbb N\\
&\hspace{12mm}\vdash\PlaceTermSeq{b}{a}\colon\mathbb N\to\mathbb N.
\end{aligned}}
Die normalisierten Ziffernfolgen werden definiert; ein allgemeiner
Existenz- und Eindeutigkeitssatz für jede Basis ist hier noch nicht
als bewiesenes Hauptresultat ausgewiesen.
\UebQuellen{\FormulaRefAuto{\BaseOK{b}\dsep a\colon\mathbb{N}\to\mathbb{N}\vdash \PlaceTermSeq{b}{a}\colon\mathbb{N}\to\mathbb{N}}[thm]}\\
\end{BandUebersicht}
